\chapter{Conclusiones y continuidad}\label{cap:conclusiones}

Este capítulo cierra la memoria. Valora el grado de cumplimiento de los objetivos fijados en
el Capítulo~\ref{cap:introduccion}, resume las aportaciones del trabajo, declara sus
limitaciones sin suavizarlas y esboza las líneas de continuidad que el sistema deja abiertas.

\section{Grado de cumplimiento de los objetivos}\label{sec:conc-objetivos}

El objetivo general, construir una plataforma de catalogación que funcionara también como
banco de pruebas reproducible de recomendación, se cumple en su totalidad. El catálogo tiene
identidad canónica propia, procedencia por obra y licencia documentada, a escala real sobre
el catálogo importado de IGDB. Una persona usuaria puede mantener su lista de pendientes,
valorar juegos e inventariar copias físicas y digitales con su formato y su edición, con
privacidad de perfil y relaciones sociales gobernadas por completo en el backend. El protocolo
de evaluación offline quedó congelado, versionado y probado antes de ejecutar ninguna
comparación. Esa comparación se llevó a cabo: las tres familias de recomendación
(basada en contenido, colaborativa e híbrida), descritas con su formulación completa en la
Sección~\ref{cap:algoritmos}, se compararon bajo ese protocolo en el
Capítulo~\ref{cap:resultados} frente a los baselines de referencia. De los objetivos específicos solo quedan pendientes, de forma parcial, dos piezas del
contrato de evaluación, que recoge la Sección~\ref{sec:conc-limitaciones}. Cada incremento del
sistema se cerró con sus criterios de éxito verificados de forma independiente, según la
disciplina descrita en el Capítulo~\ref{cap:gestion}.

\section{Limitaciones y trabajo futuro}\label{sec:conc-limitaciones}\label{sec:conc-trabajo-futuro}

Las limitaciones conocidas se declaran sin ocultarlas. La evidencia experimental publicada
procede de una única ejecución del test sobre usuarios sintéticos, setenta y nueve de ellos
evaluables, con una sola semilla: es suficiente para comparar algoritmos entre sí bajo un
protocolo común, con significación estadística ya calculada sobre esa muestra (bootstrap BCa,
Friedman, Wilcoxon y corrección de Holm), pero no permite afirmar que esa significación sea
estable frente a otra semilla o a otra población sintética sin repetir el experimento. La
cobertura de valoraciones externas del catálogo es parcial, como corresponde a una fuente
agregada de terceros y no a un panel de crítica propio; esa parcialidad se propaga, de forma
documentada, a cualquier señal de calidad externa que la use. La captura de entorno de
ejecución y consumo de recursos, ya disponible en el ejecutor paralelo para comparaciones
futuras, no se retroproyectó sobre la ejecución v15 ya cerrada, de modo que esa ejecución
concreta carece de ese detalle. El despliegue público, descrito en la
Sección~\ref{cap:despliegue}, sirve todavía el catálogo curado inicial y no el catálogo a
escala real: esa conexión se pospuso de forma deliberada al final del proyecto; la
demostración ante el tribunal se ejecuta desde el entorno local reproducible.

El banco de pruebas que esta memoria deja construido y ya explotado habilita líneas de
continuación concretas, no genéricas. La primera es repetir el protocolo publicado con nuevas
semillas y con cohortes sintéticas adicionales, para separar lo que es una diferencia
consistente entre algoritmos de lo que podría ser ruido de una única muestra, y aplicar sobre la nueva ejecución la captura de entorno y consumo de recursos que el
ejecutor paralelo ya sabe registrar. La segunda es estudiar la sensibilidad del sistema a sus propios
hiperparámetros locales (los pesos de familia de facetas, la mezcla híbrida, el factor de
diversidad), en lugar de tratarlos como constantes fijadas de una vez. La tercera, ya fuera
del núcleo de recomendación, es conectar el despliegue público al catálogo a escala real.
Cualquier extensión de este trabajo deberá crear su propio protocolo y su propio artefacto
versionado, sin alterar de forma retroactiva la evidencia ya publicada.

\section{Aportaciones y valoración final}\label{sec:conc-aportaciones}\label{sec:conc-cierre}

Las tres aportaciones anunciadas en la Sección~\ref{sec:contexto-problema} quedan entregadas.
El producto está construido, verificado y documentado. El banco de pruebas está explotado: un
corpus citable, una población sintética documentada y un protocolo offline congelado sostienen
una comparación real entre dieciséis algoritmos, con sus artefactos, semillas y
configuraciones conservados para que un tercero pueda reproducir cada cifra publicada. Y la
metodología de ingeniería que gobernó el proyecto, versionada en el repositorio con su árbol
de planificación, sus registros de decisión y su registro de agentes, permite reconstruir el
rastro completo de decisiones. Así, el trabajo entrega las dos mitades de la idea de la que
parte SavePoint: un producto de catalogación y un banco de pruebas de recomendadores cuyos
resultados se pueden volver a obtener a partir de los mismos artefactos versionados.
